\documentclass[12pt]{article}

\usepackage[spanish]{babel}
\usepackage[utf8]{inputenc}
\usepackage[T1]{fontenc}
\usepackage{mathpazo}
\usepackage{amsmath,amssymb}
\usepackage{geometry}
\geometry{margin=2.5cm}

\title{Demostración del Teorema de Taylor para una Función Bivariante}
\author{Diego Álvarez \\[4pt]
  \small \textit{Algoritmos de Optimización — INFO1159} \\[2pt]
  \small \textit{Programación No Lineal}}
\date{}

\begin{document}
\maketitle

\section{Enunciado}

Dada la función
\[
f(x_1, x_2) = 1.20x_1 + 1.16x_2 - \frac{1}{2}\bigl(2x_1^2 + x_2^2 + (x_1 + x_2)^2\bigr)
\]
con el punto \(x^T = [0.23,\; 0.41]\), el incremento \(\Delta x^T = [0.01,\; 0.01]\) y
\(t = 0.01\), demuestre el teorema de Taylor de segundo orden.

\section{Teorema de Taylor (Segundo Orden)}

Sea \(f : \mathbb{R}^n \to \mathbb{R}\) dos veces diferenciable en un entorno de \(x\).
Entonces existe \(t \in [0, 1]\) tal que
\[
f(x + \Delta x) = f(x) + \Delta x^T \nabla f(x) + \frac{1}{2} \Delta x^T \nabla^2 f(z) \Delta x,
\]
donde \(z = x + t\,\Delta x\).

\section{Solución}

\subsection{Simplificación de la función}

\begin{align*}
f(x_1, x_2)
&= 1.20x_1 + 1.16x_2 - \frac{1}{2}\bigl(2x_1^2 + x_2^2 + x_1^2 + 2x_1x_2 + x_2^2\bigr) \\
&= 1.20x_1 + 1.16x_2 - \frac{1}{2}\bigl(3x_1^2 + 2x_2^2 + 2x_1x_2\bigr) \\
&= 1.20x_1 + 1.16x_2 - \frac{3}{2}x_1^2 - x_2^2 - x_1x_2.
\end{align*}

\subsection{Cálculo del gradiente}

\[
\nabla f(x) =
\begin{bmatrix}
\frac{\partial f}{\partial x_1} \\[4pt]
\frac{\partial f}{\partial x_2}
\end{bmatrix}
=
\begin{bmatrix}
1.20 - 3x_1 - x_2 \\
1.16 - x_1 - 2x_2
\end{bmatrix}.
\]

Evaluado en \(x = [0.23,\; 0.41]\):
\[
\nabla f(x) =
\begin{bmatrix}
1.20 - 3(0.23) - 0.41 \\
1.16 - 0.23 - 2(0.41)
\end{bmatrix}
=
\begin{bmatrix}
0.10 \\
0.11
\end{bmatrix}.
\]

\subsection{Cálculo de la matriz Hessiana}

\[
\nabla^2 f(x) =
\begin{bmatrix}
\frac{\partial^2 f}{\partial x_1^2} & \frac{\partial^2 f}{\partial x_1\partial x_2} \\[4pt]
\frac{\partial^2 f}{\partial x_2\partial x_1} & \frac{\partial^2 f}{\partial x_2^2}
\end{bmatrix}
=
\begin{bmatrix}
-3 & -1 \\
-1 & -2
\end{bmatrix}.
\]

La matriz Hessiana es constante (no depende de \(x\)), por lo tanto
\(\nabla^2 f(z) = \nabla^2 f(x)\) para cualquier \(z\).

\subsection{Valor de la función en \(x\)}

\[
f(x) = 1.20(0.23) + 1.16(0.41) - \frac{3}{2}(0.23)^2 - (0.41)^2 - (0.23)(0.41).
\]

\begin{align*}
f(x) &= 0.2760 + 0.4756 - \frac{3}{2}(0.0529) - 0.1681 - 0.0943 \\
     &= 0.7516 - 0.07935 - 0.1681 - 0.0943 \\
     &= 0.40985.
\end{align*}

\subsection{Término lineal}

\[
\Delta x^T \nabla f(x) =
\begin{bmatrix} 0.01 & 0.01 \end{bmatrix}
\begin{bmatrix} 0.10 \\ 0.11 \end{bmatrix}
= 0.01(0.10) + 0.01(0.11) = 0.0021.
\]

\subsection{Término cuadrático}

Primero se calcula el punto \(z = x + t\,\Delta x\) con \(t = 0.01\):
\[
z =
\begin{bmatrix} 0.23 \\ 0.41 \end{bmatrix}
+ 0.01
\begin{bmatrix} 0.01 \\ 0.01 \end{bmatrix}
=
\begin{bmatrix} 0.2301 \\ 0.4101 \end{bmatrix}.
\]

Dado que la Hessiana es constante,
\[
\frac{1}{2} \Delta x^T \nabla^2 f(z) \Delta x
= \frac{1}{2}
\begin{bmatrix} 0.01 & 0.01 \end{bmatrix}
\begin{bmatrix}
-3 & -1 \\
-1 & -2
\end{bmatrix}
\begin{bmatrix} 0.01 \\ 0.01 \end{bmatrix}.
\]

\[
\nabla^2 f(z) \Delta x =
\begin{bmatrix}
-3(0.01) - 1(0.01) \\
-1(0.01) - 2(0.01)
\end{bmatrix}
=
\begin{bmatrix}
-0.04 \\
-0.03
\end{bmatrix}.
\]

\[
\Delta x^T \nabla^2 f(z) \Delta x =
\begin{bmatrix} 0.01 & 0.01 \end{bmatrix}
\begin{bmatrix} -0.04 \\ -0.03 \end{bmatrix}
= -0.0004 - 0.0003 = -0.0007.
\]

\[
\frac{1}{2} \Delta x^T \nabla^2 f(z) \Delta x = \frac{-0.0007}{2} = -0.00035.
\]

\subsection{Aproximación de Taylor}

\[
f(x + \Delta x) \approx f(x) + \Delta x^T \nabla f(x) + \frac{1}{2} \Delta x^T \nabla^2 f(z) \Delta x
\]

\[
f(x + \Delta x) \approx 0.40985 + 0.0021 - 0.00035 = 0.4116.
\]

\subsection{Verificación directa}

Evaluando directamente \(f\) en \(x + \Delta x = [0.24,\; 0.42]\):

\begin{align*}
f(x + \Delta x) &= 1.20(0.24) + 1.16(0.42) - \frac{3}{2}(0.24)^2 - (0.42)^2 - (0.24)(0.42) \\
&= 0.2880 + 0.4872 - \frac{3}{2}(0.0576) - 0.1764 - 0.1008 \\
&= 0.7752 - 0.0864 - 0.1764 - 0.1008 \\
&= 0.4116.
\end{align*}

\subsection{Conclusión}

\[
f(x + \Delta x) = 0.4116 = f(x) + \Delta x^T \nabla f(x) + \frac{1}{2} \Delta x^T \nabla^2 f(z) \Delta x,
\]

lo que verifica el teorema de Taylor de segundo orden. La igualdad es exacta porque
\(f\) es una función cuadrática y su matriz Hessiana es constante; por lo tanto,
el polinomio de Taylor de segundo orden coincide exactamente con la función.

\end{document}
